\documentclass[11pt]{article}
\usepackage[margin=1in]{geometry}
\usepackage[T1]{fontenc}
\usepackage[utf8]{inputenc}
\usepackage{microtype}
\usepackage{booktabs}
\usepackage{tabularx}
\usepackage{array}
\usepackage[hidelinks]{hyperref}
\setlength{\parindent}{1.2em}
\setlength{\parskip}{0.25em}
\setlength{\emergencystretch}{2em}
\newcolumntype{Y}{>{\raggedright\arraybackslash}X}

\title{Evaluation and Measurement of Runtime Safety-Monitor Cascades:\\
A Development-Only Pilot with a Corrected Claim Boundary}
\author{Tivdzua Lubem Noah\\Innopolis University}
\date{}

\begin{document}
\maketitle

\begin{abstract}
This report re-states a development-only study of selective runtime safety-monitor cascades after four follow-up analyses corrected the interpretation of the original pilot. The pilot covers three monitor stacks and three data sources and examines false-positive-rate (FPR) treatment, grouping, label choice, deployable routing, and direct end-to-end (E2E) latency. The earlier 19/75 grouping and 14/75 label changes are not isolated measurement effects: with policies, rows, predictions, thresholds, and routes frozen, the corresponding pass/fail changes are 0/75 and 1/75; retraining/reselection changes 6/75 and 4/75; the original 19/75 and 14/75 remain full-protocol contrasts. Numerical mismatches are likewise a boundary and finite-precision diagnostic rather than evidence of a general safety failure. Float64 CPU arithmetic is effectively invariant, float32 changes some decisions near frozen thresholds, and controlled CPU-versus-Tesla-T4 float32 testing produces 0/5 route and 0/5 prediction flips on the controlled mismatch cases. A prespecified dead-band analysis reports 6.12\% route ambiguity at $\epsilon=10^{-6}$, but the contribution of exact threshold ties and discrete cheap-monitor outputs to that fraction has not yet been decomposed. A small safety-availability/control-plane kill study produces no dead-band-scale continuation signal across the three stacks and does not support a standalone security-paper direction. The current evidence is therefore best treated as a pilot. Stronger claims require genuinely fresh source- and time-separated data, independent multi-rater labels, genuinely different monitor families, and a preregistered test of risk-certificate transport without retuning on confirmation data.
\end{abstract}

\section{Scope and corrected provenance}
The original project asked whether inexpensive monitors could selectively invoke a stronger monitor while maintaining harmful-response recall, controlling false positives, and reducing monitoring cost. The present evidence does not support a router-superiority claim. It is used instead as a development-only case study for evaluation and measurement.

The authorized development view contains 1,687 prompt-response examples from JailbreakBench judge-comparison, WildGuardTest, and XSTest-safe. The monitor stacks are rule $\rightarrow$ compact, compact $\rightarrow$ Qwen3Guard, and rule + compact $\rightarrow$ Qwen3Guard. The primary pilot uses audited harmful-response labels and dependency-aware grouping.

A provenance correction is essential. The legacy \texttt{final\_test} and \texttt{held\_out\_shift} splits were historically evaluated and were included in the label audit. They were excluded from the later development-only pilot, but they are not fresh confirmatory data and are not eligible for future confirmation. No result in this report is presented as fresh confirmation.

Historical v2 limitations remain part of the claim boundary. Reported v2 costs were summed measured component latencies rather than direct wall-clock E2E policy latency; historical grouping was effectively one group per example; the 35 s quantity used in bounded-cost analysis was not a demonstrated mechanically enforced runtime bound; the 1\% cost-equivalence margin was an engineering choice rather than an externally justified confirmatory margin; and fresh joint FPR+cost certification, fresh source/time-separated confirmation, and multi-rater labels were not available.

\section{Pilot evaluation results}
\subsection{Effects that remain directly attributable to the evaluation procedure}
Several pilot findings remain useful without attributing them to router quality. Empirical FPR screening and the frozen finite-sample row-and-dependency risk gate disagree on 36 of 75 deployable policy evaluations. Offline exact-ranked acquisition and the reusable deployable threshold disagree on 4 of 45 matched risk conclusions, and realized acquisition differs in all 45 matched comparisons. Direct E2E latency is slightly larger than the same-run component sum for every timed policy, but none of the 105 pairwise mean-cost orderings reverses.

These findings are descriptive properties of this development-only pilot. They are not evidence that the same effects will replicate across different monitor families or fresh data.

\subsection{Factorial decomposition of grouping and labels}
The earlier 19/75 grouping and 14/75 label results were full-protocol comparisons. Changing grouping or labels also changed folds, training data, fitted models, selected policies, thresholds, and holdouts. They therefore cannot be interpreted as isolated measurement effects.

The completed 2-by-2 decomposition separates three layers. The primary grouping contrast is dependency-aware versus singleton grouping at audited labels; the primary label contrast is audited versus original labels under dependency-aware grouping.

\begin{table}[h]
\centering
\caption{Primary grouping and label contrasts after decomposition. Counts are descriptive across 75 policy instances; the five seeds are not pooled as independent inferential observations.}
\begin{tabular}{lcc}
\toprule
Layer & Grouping pass/fail changes & Label pass/fail changes \\
\midrule
Fixed-policy measurement & 0/75 & 1/75 \\
Retraining/reselection, fixed outer holdout & 6/75 & 4/75 \\
Full protocol & 19/75 & 14/75 \\
\bottomrule
\end{tabular}
\end{table}

At the fixed-policy layer, evaluation rows, fitted models, thresholds, routes, and predictions are frozen and only the evaluation label or risk grouping changes. The small 0/75 and 1/75 effects show that the original 19/75 and 14/75 should not be described as pure grouping or label measurement effects. The larger retraining/reselection and full-protocol differences show that much of the observed variation arises through the development pipeline itself.

\section{Numerical route stability}
The archived mismatch taxonomy contains two route-threshold mismatches and five prediction mismatches. Two of the prediction mismatches are downstream consequences of route changes; three are pure decision-threshold mismatches.

The follow-up analysis separates runtime/implementation, precision, and hardware. Float64 Python, NumPy, and PyTorch CPU arithmetic is effectively invariant, with maximum perturbation $2.22\times10^{-16}$. Float32 is more sensitive near the frozen boundaries: diagnostic NumPy recomputation produces 2 route and 4 prediction flips, and PyTorch CPU produces 2 route and 8 prediction flips across the archived population. In the controlled hardware comparison, the same float32 PyTorch arithmetic produces 0/5 route and 0/5 prediction flips between CPU and Tesla T4; the compact monitor's maximum CPU-versus-T4 float32 score difference is $7.36\times10^{-9}$. T4 float16 is exploratory rather than supported and changes 2/5 routes and 5/5 predictions relative to T4 float32.

The dead-band analysis uses a prespecified grid rather than choosing a tolerance to erase the observed mismatches. The conservative non-exploratory perturbation envelope is $1.104\times10^{-7}$, so the first prespecified grid point that covers it is $\epsilon=10^{-6}$. At that point, 6.12\% of archived selective-route evaluations and 0.275\% of archived decision evaluations are classified as numerically ambiguous.

This result is a boundary/finite-precision diagnostic, not evidence of a general safety failure or universal hardware instability. In particular, the 6.12\% route-ambiguity fraction must not be interpreted as 6.12\% of routes being changed by floating-point noise. Exact threshold ties and discrete cheap-monitor outputs may contribute materially to this mass. Their contribution remains to be quantified by the planned threshold-tie diagnostic.

\section{Safety-availability/control-plane kill study}
A small development-only kill study tests both routing/bypass pressure and load/escalation pressure using the existing three stacks. It compares a fail-open budget cap, fail-closed operation without admission control, and fail-closed operation with explicit defer/reject.

At the frozen $10^{-6}$ dead-band scale, the prespecified incremental bypass continuation criterion is not met in any of the three stacks (0/3). Under the deterministic capacity model, if all accepted requests require an expensive stage with incremental cost $c_e$, serving $N$ accepted requests requires budget $B\geq N c_e$. Likewise, if escalated arrival $q\lambda$ exceeds expensive-stage service capacity $\mu$, accept-all fail-closed operation has positive backlog drift. In that overloaded regime, a hard resource bound requires either provisioning for the worst case or an explicit admission/defer/reject action rather than silent fail-open behavior.

This is a generic capacity/admission-control condition. The kill study does not establish a novel security theorem, a realizable attack result, or sufficient evidence for a standalone security-paper direction.

\section{Corrected claim boundary}
The existing evidence supports the following narrow statements:
\begin{itemize}
    \item empirical FPR screening can materially disagree with the frozen finite-sample risk gate in this pilot;
    \item ranked acquisition and a deployable threshold are operationally different in this pilot;
    \item direct E2E latency is the appropriate cost estimand, although correcting component-sum cost does not reverse mean-cost ordering here;
    \item the earlier grouping and label counts are mostly development-protocol effects rather than isolated fixed-policy measurement effects;
    \item the observed numerical mismatches are concentrated at frozen boundaries and are sensitive to precision, while controlled CPU-T4 float32 testing shows no policy flips in the five controlled mismatch cases;
    \item the control-plane kill study does not support a standalone security-paper pivot.
\end{itemize}

The evidence does \textbf{not} establish router superiority, Pareto dominance, fresh joint FPR+cost certification, fresh source/time-separated confirmation, multi-rater confirmation, universal hardware invariance, arbitrary distribution-shift robustness, production readiness, or literature novelty for the control-plane result.

The current evaluation/measurement story should therefore be treated as a pilot rather than as a completed general result. A narrow evaluation/measurement paper would require the main phenomena to replicate on genuinely different monitor families and genuinely fresh data.

\section{Next-study boundary}
The stronger research question motivated by the pilot is whether finite-sample safety/risk certificates transport across source or attack-family shift, particularly when examples are dependent or template-generated. The next work is deliberately ordered and limited.

First, one cheap diagnostic will decompose the reported 6.12\% route dead-band into exact threshold ties, discrete cheap-monitor output mass, and genuinely near-boundary floating-point perturbation. After that diagnostic, the existing data will no longer be used for discovery.

Second, before any new data are collected or evaluated, a fresh-data protocol will be frozen. It will define genuinely fresh source- and time-separated examples, independent multi-rater labels, genuinely different monitor families, dependency/template handling, and a sample-size/power calculation for a 5\% FPR certificate. The confirmatory experiment will test whether a certificate calibrated in one source or attack family remains valid under another, with no router retuning on confirmation data.

No transport result is claimed in this report because that preregistered study has not yet been run.

\section{Reproducibility and evidence map}
The numerical values in this report are tied to committed repository artifacts. The primary evidence files are:
\begin{itemize}
    \item Legacy confirmatory status: \path{data/metadata/confirmatory_split_provenance.json}
    \item Factorial decomposition: \path{reports/factorial_measurement_decomposition_v1/decomposition_summary.json}
    \item Numerical stability: \path{reports/numerical_route_stability_v1/final/summary.json}
    \item Control-plane kill study: \path{reports/safety_availability_control_plane_kill_study_v1/summary.json}
    \item Pilot timing/evaluation: \path{reports/evaluation_measurement_pilot_v1/}
\end{itemize}

\section{Conclusion}
The follow-up analyses narrow rather than expand the original claims. The strongest pilot effect remains the difference between empirical FPR screening and finite-sample risk treatment. By contrast, the earlier grouping and label counts largely reflect retraining/reselection and full-protocol changes, and the numerical mismatches are boundary-sensitive finite-precision behavior rather than a general hardware or safety failure. The small control-plane study provides a capacity/admission-control condition but no standalone security signal.

The appropriate role of the current repository is therefore as a carefully documented pilot and measurement foundation. Any stronger publication claim now requires a frozen fresh-data study of risk-certificate transport, with independent labels and no retuning on confirmation data.

\end{document}
